\subsubsection{ファンアウト1における局所群の完全な境界}
\label{subsec:fanout-one-boundary}
有限半群の擬多様体$\mathbf V$について，ファンアウト$f$で
\emph{多項式特徴データを許す}とは，任意の有限非空アルファベット
$\Sigma$，ランク上限$r_{\max}\ge1$，
および全射型付け
$s:\Sigma^+\twoheadrightarrow\mathsf S$（$\mathsf S\in\mathbf V$）
について，$s,r_{\max}$に依存してよい多項式$p_{s,r_{\max}}$が
存在し，各$L\in\Cmcf{f}{s}$と
$L(G)=L\cap\Sigma^+$を満たす縮約済み良形ファンアウト高々$f$・
ランク高々$r_{\max}$のすべての$G$に対し，
$\mathcal B^{\le r_{\max}}_{f,s}$の特徴標本$K$で
\[
 \|K\|_+\le p_{s,r_{\max}}(\|G\|,\bar\theta_G)
\]
を満たすものが存在することをいう．

\begin{theorem}[ファンアウト1の擬多様体二分法]
\label{thm:fanout-one-dichotomy}
任意の有限半群の擬多様体$\mathbf V$について次は同値である．
\begin{enumerate}[label=\textup{(\roman*)},leftmargin=*]
\item $\mathbf V$はファンアウト1で多項式特徴データを許す．
\item $\mathbf V\subseteq\mathbf{LG}$．
\end{enumerate}
さらに$\mathbf V\not\subseteq\mathbf{LG}$なら，
2元半束$U_1\in\mathbf V$への型付け$s_{\mathrm U}$に対し，
$\|G_n\|=O(n)$，$\theta_{G_n}\le4n$の縮約済み良形
ランク2表示$G_n$を持つ有限標的
$L_n\in\Cmcf{1}{s_{\mathrm U}}$が存在し，
同じ標的クラスのすべてに特徴標本を持つ任意の集合駆動型作用素は，
$L_n$の各特徴標本に長さ$3\cdot2^n-2$の語を必要とする．
\end{theorem}
\begin{proof}
(ii)$\Rightarrow$(i)．$f=1$として
定理~\ref{thm:local-group-polydata}を適用すれば
\[
 \|K\|_+=O_{s,r_{\max},\Sigma}
   (\|G\|^3\bar\theta_G)
\]
が成立する．

(i)$\Rightarrow$(ii)は対偶を示す．
$\mathsf S\in\mathbf V\setminus\mathbf{LG}$を取ると，
定理~\ref{thm:lg-characterization}(ii)から
$\mathsf S$は$U_1$と同型な部分半群を持つ．
擬多様体は部分半群で閉じるため$U_1\in\mathbf V$である．
付録~\ref{app:fanout-one-lower}の倍増言語族は
$L_n,L_n^-\in\Cmcf{1}{s_{\mathrm U}}$を満たす．
双方有限なので，それぞれ良形ランク2表示が存在し，
定理~\ref{thm:exact-reconstruction}により
$\mathcal B^{\le2}_{1,s_{\mathrm U}}$に特徴標本が存在する．
ところが補題~\ref{lem:distinguishing-pair}により
$L_n$の任意の特徴標本は唯一の差分語$w_n^*$を含み，
\[
 \|K\|_+\ge|w_n^*|+1=3\cdot2^n-1.
\]
$\|G_n\|=O(n)$，$\bar\theta_{G_n}\le4n$なので
多項式上界に反する．区別対の補題は作用素非依存であり，
末尾の主張も得られる．
\end{proof}
\begin{remark}[二分法の量化範囲]
\label{rem:fanout-one-scope}
これは\emph{擬多様体}の境界であって，
$\mathbf{LG}$の外にある個々の型付けすべてが
超多項式下界をもつという意味ではない．
下界は元の表示の厚みに対する特徴標本長についての主張であり，
Gold学習器の収束段階についての主張ではない．
ファンアウト2以上で同じ必要条件が成り立つかは未解決である．
\end{remark}
